\section{Extraction of the extension space}
\label{sec:extensions}

Set $b=b_*$ and $R_b=k[t]/(t^b)$.  We use the admissible cohomology
of \cref{def:admissible}.  The aim is to derive from
\cref{eq:reduced-hecke-congruence} an admissible Greenberg extension space
of length at least $b$.  This will contradict \cref{prop:green-length}.
The construction uses an algebra represented by $4\times4$ matrices, with block sizes
$2,1,1$ throughout.

\subsection{A finite limiting matrix algebra}

For each $i$, take the product of the Galois representations of
\cref{prop:galois-package} over all characteristic-zero systems of
$T_i$, and denote it by $\rho_i$.  Choose stable integral lattices in
the finitely many factors.  Put $B_i=T_i[\rho_i]$.  This is a
submodule of a finite product of integral matrix lattices, and hence
is finite over $\cO$ and over $T_i$.  Chebotarev approximation, as in
\cref{subsec:binary-projections}, shows that the trace and all
characteristic coefficients of every group element belong to $T_i$.

We need a bound on the number of $T_i$-module generators of $B_i$,
even though their ranks over $\cO$ may grow.  After base change, the
images of these generators under an algebra map into $M_4(R_b)$
will have one common denominator.  Expressing group matrices as
linear combinations of the generators will then give bounded
finite-stage representatives, as proved in
\cref{lem:finite-value-admissibility}.

\begin{lemma}[Uniformly many module generators]
\label{lem:bounded-algebra-generators}
There is an integer $N$ independent of $i$ such that the $T_i$-module
$B_i$ is generated by $N$ monomials in images of Galois elements.
\end{lemma}

\begin{proof}
The residual characteristic polynomials are fixed by the comparison
system.  After a fixed finite extension of $L$ killing their
semisimplification, the residual image in every factor is unipotent,
and the full integral image is pro-$p$.  All factors are unramified
outside the fixed set $S$ and the bounded number of places over
$q_i$.  The maximal pro-$p$ Galois group of this extension with such
ramification has a bounded number of topological generators: its
minimal number is the dimension of $H^1(-,\F_p)$, which the class
field sequence bounds by a fixed contribution from $S$, the units
and class group, plus one bounded contribution per new place.
This is the usual global Galois finiteness theorem
\citep[Chapter X]{NSW2008}.  Add fixed coset representatives for the
finite extension.  It follows that uniformly many elements generate
the matrix image topologically.

Their algebraic span is already closed, because any $\cO$-submodule
of the finite integral matrix lattice is finite and closed.  Inverses
are polynomial expressions with integral coefficients by the
characteristic equation and the unit determinant.  Hence the same
bounded list generates $B_i$ as an algebra.

Reduce modulo the maximal ideal of $T_i$.  Matrix algebras of size
four satisfy the degree-eight standard polynomial identity
\citep{AmitsurLevitzki1950}, and this identity passes to the quotient.
Shirshov's height theorem bounds the number of power factors needed
to span the quotient, with base words of length less than eight;
the bound depends only on the identity and the number of generators
\citep[Section~2, Theorem~1 and its proof]{Shirshov1957}.
Each base word before reduction
is a group image and satisfies its degree-four characteristic
polynomial over $T_i$.  Reducing its powers to exponents at most
three leaves a uniformly bounded list of monomials.  These span
$B_i$ modulo the maximal ideal of $T_i$.  Since $B_i$ is finite over $T_i$,
Nakayama's lemma proves the assertion.
\end{proof}

Choose the nonprincipal ultrafilter $\mathcal U$ used in
\cref{sec:tame}, and form the algebraic ultraproduct
$A=\prod_{\mathcal U}T_i$.  The constant coefficients of
\cref{eq:reduced-hecke-congruence} have compact $p$-adic ultralimits,
giving a map $A\to\cO\subset k$; let $\mathfrak q$ be its kernel.
Taking the ultralimit of every coefficient of $t$ gives a map
$A\to R_b$ with this constant coefficient.  Every element outside
$\mathfrak q$ has unit image in $R_b$, so this map extends to
$A_{\mathfrak q}$.  Since $R_b$ is henselian, it extends uniquely
to the henselization.  Put
\begin{equation}
 T=(A_{\mathfrak q})^h,\qquad
 \lambda:T\longrightarrow R_b,\qquad
 B=\left(\prod_{\mathcal U}B_i\right)\otimes_A T.
 \label{eq:limiting-algebras}
\end{equation}
The residue field of $T$ is $k$: the image before localization
contains $\cO$, and localizing at the kernel inverts every nonzero
constant.
Write $\mathfrak m=\ker(T\to k)$.

There is a useful faithful matrix realization.  Put
\begin{equation}
 C=\prod_{\mathcal U}\prod_{\pi\text{ at }i}\overline{\Q}_p,
 \qquad \mathcal R=C\otimes_A T.
 \label{eq:ambient-ring}
\end{equation}
The injections $A\hookrightarrow C$ and
$\prod_{\mathcal U}B_i\hookrightarrow M_4(C)$ remain injective after
tensoring with the flat localization and henselization $T$.
Thus $T\hookrightarrow\mathcal R$ and
$B\hookrightarrow M_4(\mathcal R)$ are faithful.  The ring $C$ is
absolutely flat.  Its residue fields are algebraically closed
ultraproducts obtained from ultrafilters on the total set of systems
which prolong $\mathcal U$.  Base change to $T$ is a filtered
combination of localization and etale operations.  These preserve
absolute flatness and reducedness, so $\mathcal R$ is again
absolutely flat and reduced.  Its field tests can be made in these
algebraically closed ultraproducts; finite etale choices split over
them.  This observation will allow finite matrix identities to be
tested on actual arrays of characteristic-zero systems.

Let $\Gamma$ denote the ultraproduct of the Galois groups unramified
outside $Sq_i$, with the prescribed local subgroups.  The bounded
monomials in \cref{lem:bounded-algebra-generators} still generate $B$
after both base changes.  The monomials may vary with $i$; their number
is fixed.  They are products of group images, so
$B$ is generated by $\rho_*(\Gamma)$, and its trace is $T$-valued.
The limiting trace under $\lambda$ is
\begin{equation}
  \tr(f\oplus x\oplus y),\qquad
  f=V(-2)\psi^{\pm1},\quad x=\eps^{-1},\quad y=\eps^{-2}.
 \label{eq:limiting-blocks}
\end{equation}
Here the sign follows the reciprocity choice in
\cref{sec:conventions}, and either sign is permitted in
\cref{prop:green-length}.  Indeed, the split good polynomials are
those of the theta comparison.  The factor $\cC^m$ tends uniformly
to one on the fixed torsion branch, and Frobenius approximation at
each finite stage may be made to arbitrary increasing precision.
This proves \cref{eq:limiting-blocks} on all sequences of elements.

The corresponding diagonal representation factors through $B$.
Here reduction modulo $\mathfrak m$ has coefficient field $k$, of
characteristic zero.  Its diagonal blocks are absolutely irreducible and
pairwise distinct: $V|_{G_L}$ is absolutely irreducible by the choice
of $L$ in \cref{lem:auxiliary-fields}.  Their group images span
$M_2(k)\times k\times k$, and hence, by Nakayama, span
$M_2(R_b)\times R_b\times R_b$.  If a relation vanishes in $B$,
its product with every group-algebra element has trace zero.  The
trace identity \cref{eq:limiting-blocks} and the perfect trace
pairing on this matrix product show that its image in the diagonal
representation is zero as well.

\subsection{Corners and admissible representations}

We spell out the elementary matrix-algebra construction to keep
track of finiteness and of the coefficient ring; compare the general
framework of \citet[Sections~1.3--1.5]{BellaicheChenevier2009}.
Choose an element of $B$ acting with four distinct prescribed
eigenvalues in the residual diagonal representation.  Its
characteristic coefficients belong to $T$: Newton identities express
them through traces, and the integers $1,2,3,4$ are units in this
local ring with characteristic-zero residue field.  Henselianity splits its characteristic
polynomial into four linear factors with pairwise unit differences.
The associated polynomial projectors are orthogonal idempotents
$e_1,e_2,e_x,e_y$ summing to one.  In every ambient field their ranks
are one, since the split characteristic polynomial has four simple
roots.  The first two lie over the $f$ block.

Lift the two off-diagonal matrix units in that residual $2\times2$
block.  Their products in the rank-one corners are unit scalars;
rescaling gives inverse matrix units between $e_1$ and $e_2$.
A diagonal corner is precisely $T e_j$: in the faithful rank-one
realization each of its elements equals its trace times $e_j$.
Using the matrix units to identify the first two positions, let
$D_{uv}$, for $u,v\in\{f,x,y\}$, be a single-entry corner from $v$
to $u$.  These finite $T$-modules have the usual composition products,
and $D_{uu}=T$.  All crossing diagonal cycles map to zero under
$\lambda$, since the diagonal representation annihilates every
crossing entry.  In particular
\begin{equation}
 \lambda(D_{uv}D_{vu})=0\quad(u\ne v).
 \label{eq:crossing-cycles}
\end{equation}

Our eventual extension module will be
$D_{fy}\otimes_{T,\lambda}R_b$.  A functional on it should retain
the extension of $y$ by $f$ and discard the other crossing entries.
To respect multiplication, it must also kill the discarded path from
$y$ through $x$ to $f$.  We will first prove
$D_{xy}=D_{xf}D_{fy}$, so that this path contains a crossing diagonal
cycle and vanishes under \cref{eq:crossing-cycles}.  We will then show
that $D_{fy}$ is nonzero in every ambient field, giving the vanishing
Fitting ideal from which the length bound follows.

\begin{lemma}[Finite-value admissibility]
\label{lem:finite-value-admissibility}
If $\Phi:B\to M_4(R_b)$ is a $T$-linear algebra homomorphism with
scalar map $\lambda$, then $\Phi\rho_*$ has admissible matrix
coefficients.  The same holds with $k$ in place of $R_b$.
Algebraic local identities holding on all local sequences have the
uniform increasing-precision meaning in \cref{def:admissible}.
\end{lemma}

\begin{proof}
Let $b_{ij}$, $1\leq j\leq N$, be the monomial generators from
\cref{lem:bounded-algebra-generators}, and $b_j$ their images in $B$.
Only the fixed finite list of matrices $C_j=\Phi(b_j)$ is needed.
One denominator clears all coefficients of all $C_j$ in the fixed
basis of $R_b$.

For any increasing precision $s_i$, the compact surjection
$T_i^N\to B_i$ induces a surjection on the finite quotients modulo
$p^{s_i}$.  Choose a section on that finite set.  Composing with
$\rho_i$ produces locally constant coefficient functions
$a_{ij}(g)\in T_i$ such that
\[
 \rho_i(g)\equiv\sum_{j=1}^N a_{ij}(g)b_{ij}
                       \pmod {p^{s_i}B_i}.
\]
Apply \cref{eq:reduced-hecke-congruence} and multiply the coefficients
by successively accurate approximations to the fixed $C_j$.
The resulting continuous matrix functions have the one denominator
already chosen, and their limit is $\Phi\rho_*$ by $T$-linearity.
The precision may be slowed to the minimum of $s_i$ and $M_i$ minus
this fixed loss; the bound $M_i\leq n_i-c_b$ already respects the
tame relation.  The variable element $(p^{s_i})_i$ has image zero
under $\lambda$ in every coefficient of $t$, so the remainders vanish.
Although the constant element $p$ is a unit in $T$, this does not
change the assertion about the variable sequence.

If a matrix identity holds for all sequences but fails uniformly to
increasing precision, choose a violating local element at each of
a filter-large set of stages.  That sequence contradicts the identity
in the limit.  The same argument applies to local splittings after
choosing the finitely many splitting coefficients.  In particular
no coordinatewise lifting of the henselian idempotents is required.
\end{proof}

The local subspaces of the characteristic-zero representations need
not themselves survive a corner quotient.  We transport their rank
conditions as polynomial identities in corner entries instead.  The
following lemma makes these identities intrinsic to $B$ and hence
preserves them under every algebra map used below.
\begin{lemma}[Corner determinant identities]
\label{lem:corner-minors}
Choose four rows of a stack of $4\times4$ matrices in $B$, with
column indices $1,2,x,y$.  If the row indices consist of the column
indices with $y$ replaced by $u$, then its signed determinant is a
well-defined element of $D_{uy}$.  The same assertion holds for a square minor whenever canceling
common row and column indices leaves precisely one surplus row
and one surplus column; if no indices remain, its value lies in $T$.
If the ordinary determinant is zero in every ambient field, this
corner element is zero.  Its identity is preserved by algebra maps
out of $B$.
\end{lemma}

\begin{proof}
Each signed determinant monomial has entries directed from its
column index to its row index.  At every vertex the incoming and
outgoing counts agree except for a path from $y$ to $u$.  Decompose
the directed graph into that path and cycles.  Compose the path in
$D_{uy}$; each cycle is a scalar in its rank-one diagonal corner and
hence in $T$.  The inverse matrix units identify the two $f$
positions whenever needed.  This gives the claimed element.  In a
rank-one matrix realization it is exactly the usual product of
entries, independent of the chosen bases, since the intermediate
basis changes cancel.  Faithfulness in $M_4(\mathcal R)$ and
reducedness of $\mathcal R$ prove the vanishing assertion.  The
construction uses sums and corner products, so is preserved by
algebra homomorphisms.
\end{proof}

\subsection{Eliminating the cyclotomic extension corner}

\begin{proposition}
\label{prop:cyclotomic-corner}
The corner modules satisfy
\begin{equation}
 D_{xy}=D_{xf}D_{fy}.
 \label{eq:xy-factorization}
\end{equation}
\end{proposition}

\begin{proof}
Let
\[
 Q_{xy}=D_{xy}/(D_{xf}D_{fy}+\mathfrak mD_{xy}).
\]
A $k$-linear functional $Q_{xy}\to k$ defines a triangular
representation, retaining the separate $f$ block and an extension
of $y$ by $x$.  Products through $f$ have been killed, and other
crossing products vanish by \cref{eq:crossing-cycles}; thus the
triangular assignment is an algebra map out of $B$.  It yields an
injection
\begin{equation}
 \Hom_k(Q_{xy},k)\hookrightarrow H^1_{\mathrm{adm}}(\Gamma,k(1)).
 \label{eq:xy-cohomology}
\end{equation}
Indeed, a split Galois extension has a splitting preserved by the
algebra generated by the group, hence by $B$.  The idempotents then
force this splitting to be the diagonal $y$ line.  Its invariance
annihilates the whole functional.  Admissibility follows from
\cref{lem:finite-value-admissibility}.

We verify the strict local conditions on every class in
\cref{eq:xy-cohomology}.  At a conjugate place $\bar w\mid p$, set
$S_g=\rho_*(g)-Y_{\bar w}(g)$ for local sequences $g$.
In every characteristic-zero system these matrices have a common
kernel by \cref{lem:two-contacts}.  Take four rows of their stack
with types $f,f,x,x$.  The four-row determinant is an element of
$D_{xy}$ by \cref{lem:corner-minors}, and it is zero by the common
kernel.  Its identity therefore holds in the triangular quotient.

On the three competing columns $f_1,f_2,x$, one may choose two
$f$ rows and one $x$ row with an invertible $3\times3$ minor.
Here the inverse matrix units of the $f$ block transport the selected
coordinate rows to the two required output positions.  Writing these
units as $E_{ab}$, the rows can be represented by
$e_1E_{1a}S_g$ and $e_2E_{2b}S_h$, with $a,b\in\{1,2\}$ allowed
to coincide.  These operations preserve the common kernel and the
corner identities of \cref{lem:corner-minors}.
For the $f$ part this follows from
\[
 \Hom_{G_{L_{\bar w}}}(y,f)
       =H^0(L_{\bar w},V)=0;
\]
there is no common kernel, so finitely many rows span the two-column
dual.  The tame character is trivial locally at these fixed places.
For the $x$ row use $x\ne y$ on the local group.
Normalize the $y$ coordinate to one and solve these three equations
for the other coordinates.  Each additional $x$ row is consistent
with this solution because its four-row determinant is zero.
The $f$ coordinates are zero, all the remaining $f$ equations
already hold, and the $y$ rows are zero.  Thus a common $y$-eigenline
exists in the triangular extension: its restriction at $\bar w$
splits.

At $w\in\Sigma_p$, use the common quotient of character $X_w$ in
\cref{lem:two-contacts} and apply the transposed minor argument.
The competing-column condition becomes
\[
 \Hom_{G_{L_w}}(f,x)=H^0(L_w,V)=0,
\]
using $V\simeq V^*(1)$, together with $x\ne y$.
The resulting quotient of character $x$ splits the same extension
of $y$ by $x$.  \Cref{eq:last-character,eq:first-character} supply the required
characters on all local elements.  This proves strictness at every
place above $p$.

Away from $p$ the extension is unramified, as follows.  Outside the
controlled places this holds for the original representations.
At a binary control place, the image of $\Delta_v$ is a unit in $T$.
The coprime factorization of the full Frobenius polynomial therefore
defines $P_v^{\mathrm{bin}}$ over $B$, and the identity
\cref{eq:binary-inertia-identity} holds there by faithful field
testing.  In the triangular representation this projection is the
identity on the $x,y$ block: both of its Frobenius roots are in $A_v$
and are separated from the $f$ roots.  Hence inertia acts trivially
on that extension.  This includes the moving places $q_i$.

At a Steinberg control place, the rank bound of
\cref{lem:steinberg-control} gives two-by-two minors of
$\rho_*(\iota)-1$ with rows $f,x$ and columns $f,y$, which are zero
in $D_{xy}$.  In the triangular representation they say
\[
  (f(\iota)-1)_{ab}\,u(\iota)=0
\]
for every matrix entry and every inertia element, where $u$ is the
extension entry.  The comparison elliptic Steinberg representation
has a nonzero actual inertia difference for some element $h$.
Thus $u(\iota)=0$ whenever $f(\iota)-1\ne0$, in particular for $h$.
If $f(\iota)=1$, apply the same assertion to $\iota h$ and use the
cocycle equation and the unramifiedness of $x,y$ to obtain
$u(\iota)=u(\iota h)-u(h)=0$.

By admissibility and unramifiedness at the moving places, the class
now descends to ordinary continuous cohomology over $L$, by
\cref{lem:admissible-specialization}.  Its coefficients are $k(1)$, it is
unramified away from $p$, and it is zero locally at all $p$-places.
Kummer theory identifies such classes with $p$-units tensored with
$k$, subject first to zero valuations at the $p$-places and then
to zero local logarithms.  The valuation condition leaves global
units.  The CM biquadratic field $L$ has unit rank one, and a
nontrivial fundamental unit has nonzero $p$-adic logarithm at a
split $p$-place: if its image in $\Q_p^\times$ had logarithm zero,
being a unit it would be a root of unity, and injectivity of the
field embedding would make the global unit torsion.  Thus the
local logarithm is injective on this one-dimensional space, and the
class vanishes.  We have shown $Q_{xy}=0$.  Nakayama's lemma for
the finite module $D_{xy}/D_{xf}D_{fy}$ proves
\cref{eq:xy-factorization}.
\end{proof}

\subsection{A nonzero corner in every characteristic-zero field}

\begin{lemma}[Excluding a stable last line]
\label{lem:no-stable-y-line}
The image of $D_{fy}$ is nonzero at every field evaluation of
$\mathcal R$.
\end{lemma}

\begin{proof}
Suppose its image vanished at such a field.  By
\cref{eq:xy-factorization}, $D_{xy}$ also vanishes there, so the
one-dimensional $e_y$-space is globally stable.
We will obtain a rank-one geometric character with conjugate weight
sum three.  The residual corner character $y=\eps^{-2}$ has
weight sum four.  Evaluating on a fixed pair of conjugate local units
will turn this discrepancy into a scalar unit that would have to
vanish in the ambient field.

We first identify the stable $e_y$-space at a conjugate place
$\bar w\mid p$ with the
first line in \cref{lem:two-contacts}.  Choose the two $f$ rows and
one $x$ row used above from the stack
$\rho_*(g)-Y_{\bar w}(g)$, and project their columns to the complement
of $e_y$.  Their $3\times3$ determinant is a scalar in $T$ and is
a unit: its reduction is the invertible minor in
$f\oplus x\oplus y$.  On a filter-large set of actual systems, the common kernel
is consequently a unique line with nonzero $e_y$ coordinate.  Once
$e_y$ is stable, all competing entries in its column are zero.
The invertible complementary minor therefore forces the existing
common kernel of the matrices $\rho_*(g)-Y_{\bar w}(g)$ to be
the $e_y$ line.  Its character is thus $Y_{\bar w}$.

Here these assertions really do refer to arrays of actual
characteristic-zero systems.  Stability need only be checked on
the uniformly bounded list of module generators in
\cref{lem:bounded-algebra-generators}.  The rank, inverse matrix-unit
relations and the competing-column condition also use finitely
many matrices.  Include also $\rho_*(g)$ and $\rho_*(g^c)$ for
the fixed unit element selected below; this only enlarges that
finite list.  Elements introduced by henselization can be
represented in the algebraically closed residue ultraproduct of
\cref{eq:ambient-ring}, with their finite relations.  Vanishing or
nonvanishing of each of these finitely many matrix expressions then
holds on a filter-large set of actual systems.  Thus at almost every
such system its semisimple global representation has a
one-dimensional constituent whose local character is the prescribed
$Y_{\bar w}$.  This passage uses finite matrix equations, not an
assertion about infinitely many coordinatewise projectors.

Denote this actual character by $\xi$.  It is de Rham above $p$,
hence Hodge--Tate with semisimple rank-one inertia, and therefore
locally algebraic
\citep[Chapter~III, Section~1.2 and Appendix~A7, Theorem~3]{Serre1968}.
Global class field theory and the Serre-group construction then
identify it with the $p$-adic realization of an algebraic Hecke
character; see \citet[Chapter~III, Sections~2.2--2.3, including
the proof, and Chapter~II, Section~2.7]{Serre1968}.
Here the construction is over the coefficient field.  The
sums of its weights at conjugate embeddings are parallel: apply its
infinity type to the global units of the real subfield.  Write
$w(\xi)\in\Z$ for this common sum.

The character occurs in a shifted copy of a cuspidal factor
$\tau[a]$ in \cref{eq:four-dimensional-construction}.  If $a>1$,
the same rank-one constituent occurs in all its consecutive shifts.
Its weight $3-h_4$ at $\bar w$ would then have an adjacent integer
weight in the full representation.  This is impossible: it is the
lowest weight there, separated from the next by $h_4-h_3$, which
was chosen arbitrarily large.  Hence $a=1$.
At every unramified place, the roots of a unitary generic cuspidal
local representation have absolute values strictly between
$Q^{-1/2}$ and $Q^{1/2}$
\citep[Theorem~7.5]{Tadic1986} and
\citep[Theorem~9.7]{Zelevinsky1980}.  After the cohomological shift in
\cref{eq:four-dimensional-compatibility}, the value of $\xi$ has
absolute size strictly between $Q$ and $Q^2$.  On the other hand,
for an algebraic Hecke character with parallel conjugate weight
sum $w(\xi)$ its absolute value is $Q^{w(\xi)/2}$: dividing out
this norm power leaves a unitary character.  Consequently
\begin{equation}
                    w(\xi)=3.
 \label{eq:rank-one-weight}
\end{equation}

There is a uniform open subgroup of the local units above $p$ on
which the smooth factors of all these characters are trivial.
Indeed the compact open level $J$ at $p$ is fixed.
Bernstein's decomposition theorem gives finitely many inertial
cuspidal-support classes among representations with $J$-invariants
\citep[Section~2.3]{BernsteinDeligneKazhdan1986}; local Langlands therefore
gives finitely many semisimple inertia types.  The smooth factor of
a rank-one constituent belongs to this finite list.  This bounds
its conductor independently of the growing algebraic weights.
Choose in this fixed open subgroup an element $g$ with
$\eps(g)\ne1$, deep enough that $\eps(g)$ is not torsion, and use
its conjugate $g^c$ at the opposite place.  \Cref{eq:rank-one-weight} implies at every actual constituent under
consideration
\begin{equation}
                    \xi(g)\xi(g^c)=\eps(g)^{-3}.
 \label{eq:weight-three-unit}
\end{equation}
Thus it implies the same identity in the evaluated ambient field.

But the scalar corner element
\[
 (e_y\rho_*(g)e_y)(e_y\rho_*(g^c)e_y)
                      -\eps(g)^{-3}e_y
\]
is a unit in $T e_y$.  Its reduction in the residual diagonal
representation is
$(\eps(g)^{-4}-\eps(g)^{-3})e_y\ne0$, since the residual $y$ is
$\eps^{-2}$ and cyclotomic is invariant under conjugation.
A unit cannot vanish in a field evaluation, contradicting
\cref{eq:weight-three-unit}.  The assumed vanishing of $D_{fy}$
is impossible.
\end{proof}

\begin{corollary}
\label{cor:fitting-vanishing}
One has
\begin{equation}
                       \Fitt_{0,T}(D_{fy})=0.
 \label{eq:fitting-vanishing}
\end{equation}
\end{corollary}

\begin{proof}
If $a\in T$ annihilates $D_{fy}$, then its image at every ambient
field is zero by \cref{lem:no-stable-y-line}.  Since $T$ embeds in
the reduced ring $\mathcal R$, this gives $a=0$.
For a finite module over an arbitrary commutative ring its zeroth
Fitting ideal is contained in its annihilator: choose finitely many
generators and apply the adjugate identity to each square matrix of
relations.  This does not require a finite presentation.  It proves
\cref{eq:fitting-vanishing}.
\end{proof}

\subsection{The Greenberg extension module and the contradiction}

\begin{proposition}[Extension extraction]
\label{prop:extension-extraction}
There is an injective $R_b$-linear map
\begin{equation}
 \Hom_{R_b}(D_{fy}\otimes_{T,\lambda}R_b,R_b)
       \hookrightarrow
 \GrH(L,V\otimes\psi^{\pm1}).
 \label{eq:green-injection}
\end{equation}
Its source has $R_b$-length at least $b$.
\end{proposition}

\begin{proof}
Put $M=D_{fy}\otimes_{T,\lambda}R_b$.  A functional $M\to R_b$
gives a triangular extension of $y$ by $f$ while retaining the
separate $x$ block.  We check the only additional product relation.
The path from $y$ through $x$ to $f$ has image in
\[
 D_{fx}D_{xy}=D_{fx}D_{xf}D_{fy}
\]
by \cref{eq:xy-factorization}.  Its first two factors form a crossing
diagonal cycle, killed by \cref{eq:crossing-cycles} after applying
$\lambda$.  All other discarded products are killed by the same
cycle relation or by the diagonal representation.  Hence the
triangular assignment is an algebra homomorphism out of $B$.
As in \cref{eq:xy-cohomology}, a split extension must have its
splitting preserved by $B$ and then by its idempotents; it forces
the functional to be zero.  This proves injectivity into extension
cohomology.  Moreover
\[
                  \Hom(y,f)=V\otimes\psi^{\pm1},
\]
and \cref{lem:finite-value-admissibility} proves admissibility with
coefficients in the full ring $R_b$.

At a conjugate $p$-place, use the same common-kernel stack
$S_g=\rho_*(g)-Y_{\bar w}(g)$.  Choose four rows of types
$f,f,x,f$, with either of the two $f$ positions as the final row.
By \cref{lem:corner-minors} the vanishing four-row determinants are
identities in the corresponding entry of $D_{fy}$, and they
survive in the triangular extension.  The first two $f$ rows and
one $x$ row give a $3\times3$ minor on the competing columns whose
residue is invertible by $H^0(L_{\bar w},V)=0$ and $x\ne y$.
It is therefore a unit over $R_b$.  Solve these equations with
$y$ coordinate one.  Varying the final $f$ row proves that every
other $f$ equation is satisfied.  The $x$ coordinate is zero and
the $y$ rows vanish.  This produces a common $y$-eigenline and
splits the local extension over the full Artin ring, including its
nilpotent coefficients.

The Greenberg condition is full at $\Sigma_p$, so no splitting is
required there; it is relaxed at the remaining fixed bad places.
Outside those places and the moving places the representations are
unramified, as required for admissible classes.  Thus the preceding
local splitting is precisely the one-sided condition of
\cref{prop:green-length}, and gives \cref{eq:green-injection}.

Fitting ideals commute with base change for finite modules, so
\cref{eq:fitting-vanishing} gives $\Fitt_{0,R_b}(M)=0$.
The finite module $M$ over the principal Artin ring $R_b$ has an
elementary-divisor decomposition
\[
              M\simeq\bigoplus_j R_b/(t^{a_j}),
                   \qquad 1\leq a_j\leq b.
\]
Its zeroth Fitting ideal is $(t^{\sum_j a_j})$, with the evident
interpretation for free summands $a_j=b$.  Vanishing is equivalent
to $\sum_j a_j\geq b$.  Each summand and its $R_b$-dual have length
$a_j$, since $\Hom_{R_b}(R_b/(t^{a_j}),R_b)=(t^{b-a_j})$.
Thus the source in \cref{eq:green-injection} has length at least $b$.
\end{proof}

We can now finish the arithmetic argument.  Outside the strict case,
$b=b_*=1$, whereas \cref{prop:green-length} gives a Greenberg module of
length zero.  In the strict case, $b=b_*=3$, whereas the same
proposition bounds its length by two.  Both contradict
\cref{prop:extension-extraction}.  Consequently the torsion
alternative for the conductor-one Heegner point in
\cref{eq:torsion-assumption} cannot occur.  The arithmetic reduction
of \cref{prop:arithmetic-reduction},
Gross--Zagier and Kolyvagin then give the analytic order, the
Mordell--Weil rank and the finiteness of the whole Tate--Shafarevich
group asserted in \cref{thm:main}.
